\documentclass[12pt]{article}
\usepackage[margin=0.75in]{geometry}
\usepackage{lmodern}
\usepackage{fancyhdr}
\usepackage{titlesec}
\usepackage{enumitem}
\usepackage{xcolor}
\usepackage[hidelinks]{hyperref}

\setlength{\parindent}{1.1em}
\setlength{\parskip}{0.4em}
\setlength{\headheight}{15pt}
\titleformat{\section}{\large\bfseries\scshape}{}{0pt}{}
\titleformat{\subsection}{\normalsize\bfseries}{}{0pt}{}
\pagestyle{fancy}
\fancyhf{}
\fancyfoot[C]{\thepage}
\renewcommand{\headrulewidth}{0.4pt}

\fancyhead[L]{Whately: Making Clear Claims}
\fancyhead[R]{Argument Lab}
\begin{document}
\begin{center}
{\LARGE\bfseries Argument Lab}\par\vspace{0.25em}
{\large\scshape Identifying and Clarifying Propositions}\par\vspace{0.5em}
\rule{0.82\textwidth}{0.4pt}
\end{center}

\section*{Part 1: Identify Quality and Quantity}
For each proposition, state whether it is affirmative or negative, and whether it is universal or particular. Then write it in standard form (A, E, I, or O).

\begin{enumerate}[leftmargin=*]
\item Every citizen has the right to vote.

\item No one under eighteen may purchase alcohol.

\item Some metals conduct electricity.

\item Not every student completed the assignment.

\item All triangles have three sides.

\item Some birds cannot fly.

\item No even number is odd.

\item Most politicians keep their promises.
\end{enumerate}

\section*{Part 2: Determine Distribution}
For each proposition, indicate whether the subject term and the predicate term are distributed or undistributed.

\begin{enumerate}[leftmargin=*]
\item All mammals are warm-blooded.

\item No reptiles are mammals.

\item Some philosophers are logicians.

\item Some philosophers are not logicians.

\item All squares are rectangles.

\item No prime numbers greater than 2 are even.
\end{enumerate}

\section*{Part 3: Translate into Standard Form}
Rewrite each of the following statements as one of the four standard forms (All S are P, No S are P, Some S are P, Some S are not P). Be prepared to explain your choices.

\begin{enumerate}[leftmargin=*]
\item Only registered students may attend the lecture.

\item None but the lonely heart can know my sadness.

\item Every student except freshmen must take the exam.

\item The only people allowed in the room are faculty.

\item Not all that glitters is gold.

\item Nothing is both a square and a circle.

\item Whenever it rains, the ground gets wet.

\item No one succeeds without effort.
\end{enumerate}

\section*{Part 4: Spot the Problem}
Each of the following statements is unclear or ambiguous in its quantity or quality. Rewrite each one so that its logical form is explicit and precise.

\begin{enumerate}[leftmargin=*]
\item “Politicians lie.” (What exactly is being claimed?)

\item “Some people are trustworthy.” (Does this mean at least one, or does it suggest not all?)

\item “Students are not prepared for college.” (Is this a claim about all students or some?)

\item “No one should be allowed to criticize the government.” (Is this universal or particular in scope?)
\end{enumerate}

\section*{Part 5: Window Note Practice}
Choose one proposition from Part 3 and create a window note that shows:
\begin{itemize}
\item The original wording
\item The standard logical form
\item The subject and predicate clearly labeled
\item Whether the statement is universal or particular, affirmative or negative
\end{itemize}

At the bottom, write a short paragraph explaining why making the quantity and quality explicit changes how the claim could be supported or challenged.

\end{document}